\documentclass[12pt]{article}

% ---------- Margins ----------
\usepackage[margin=1in]{geometry}

% ---------- Font (Closest to Times in pdfLaTeX) ----------
\usepackage{newtxtext}

% ---------- Spacing ----------
\usepackage{setspace}
\doublespacing{}

% ---------- Paragraph Formatting ----------
\setlength{\parindent}{0.5in}
\setlength{\parskip}{0pt}

\begin{document}

% ---------- MLA Heading (Flush Left) ----------
\noindent
Johnny Alejandro Rojas\\
February 20, 2026\\
ENGL 101-A\\
LaEsha Sanders, M.F.A.

\vspace{1em}

% ---------- Title ----------
\begin{center}
This I believe
\end{center}

A common struggle shared among Christians is the effort to hold faith steady
when life feels uncertain. When difficulties stretch longer than expected, it
becomes easy to wonder whether such trials are truly necessary. The human heart
often leans toward a simple conclusion: if God desires good for His people, why
permit hardship at all? From this perspective, suffering can appear like an
interruption to blessing rather than a part of it.

There is a quiet habit in human nature to define what is “good” by what feels
secure. A stable job, financial safety, peace within the family, restored
relationships, these things naturally seem to be the signs of divine favor. They
are not wrong desires. In fact, they reflect a longing for order and well-being.
Yet one may ask whether these visible goods are the whole picture. Is blessing
only what brings comfort? Or might it include something deeper that is not
immediately pleasant?

Scripture presents God not as distant, but as One who guides with a premeditated
intention. In Book of Psalms 32:8, God speaks of instructing and teaching, of
watching over with care. This suggests direction rather than randomness. If
divine guidance is real, then trials cannot be meaningless. They must serve a
purpose that may not be obvious at first glance.

Human understanding, however, is limited. The mind tends to measure life by the
moment it exists in. It sees what is missing before it sees what is being formed. It
desires the gift without always considering the weight of the gift. Yet history
and daily experience show that blessings given without preparation can harm.
Success without discipline can inflate pride. Influence without maturity can
mislead. Even love, when pursued without patience, can become destructive. In
this light, trials begin to look less like punishment and more like preparation.
They build a kind of strength that comfort alone cannot produce.

An example that often stands at the center of this reflection is the life of
Joseph. His story does not begin with comfort. It begins with betrayal. Sold by
his own brothers, carried away from his father’s house, placed in a foreign
land—none of these events resemble blessing. For a season, it would have been
reasonable to conclude that his dreams were illusions and that divine favor had
withdrawn.

Yet the narrative does not end in the pit or in the prison. Through years of
false accusation, forgotten promises, and silent waiting, something was being
formed. Joseph learned administration in Potiphar’s house. He learned endurance
in confinement, and control of his passions. He learned discernment through
suffering. None of these lessons appeared glorious at the time. Still, they
became necessary.

When famine later struck the region, it was not merely a capable ruler that
Egypt required, but a prepared one. Joseph’s earlier trials positioned him to
interpret Pharaoh’s dreams and to design a plan that would preserve grain during
years of abundance. The very suffering that seemed unnecessary became the means
by which an entire nation survived hunger. What appeared as personal misfortune
unfolded into collective preservation.

This pattern invites reflection. If Joseph had been elevated without hardship,
would he have possessed the depth required for such responsibility? If his early
dreams had been fulfilled immediately, would he have understood the weight of
leadership? The delay, painful as it was, formed the capacity to steward what
was coming.

There is also something revealing about hardship. It exposes the fragile belief
that one is fully self-sufficient. When plans fail and control slips away, the
heart becomes aware of its dependence. What was once an abstract idea of reliance
on God turns into a lived reality. In that place, one begins to see that faith
is not merely agreement with doctrine but trust formed through need of knowing
what is to come.

The Bible does not present this as a theory, but as a pattern. Growth often
comes before elevation. Waiting often comes before fulfillment. The fallen
nature of humanity, inclined toward self-centered desire, does not easily accept
delay. Yet delay may be a form of care. What appears as silence may be
instruction. What feels like denial may be protection.

The question then remains: what is truly best? Is it employment, stability,
acceptance, reconciliation? These are meaningful goods, but perhaps they are not
the highest good. Christian teaching suggests that what is best is shaped in
heaven before it appears on earth. It is defined not only by immediate
satisfaction, but by eternal purpose. The human mind cannot fully grasp that
scope. It sees in fragments; God sees the whole.

From this perspective, trials begin to reflect love rather than absence. Love
does not always grant what is desired in the moment. Sometimes it strengthens
before it gives. Sometimes it withholds to prevent harm. Sometimes it waits
until the heart is ready to carry what it has asked for.

Curiosity about suffering, then, opens the door to humility. It allows one to
admit that human judgment is partial. It invites reflection on whether hardship
might be shaping the soul for something greater than comfort. And in that
reflection, trials no longer stand only as obstacles, but as instruments quietly
building the capacity to receive what, in time, will truly bless.

\end{document}
